\documentclass[12pt]{article}
\usepackage[utf8]{inputenc}
\usepackage[T1]{fontenc}
\usepackage{hyperref}
\title{Seychelles French Creole Language Resources for K-12 Curriculum Development}
\author{Compiled by Claude Sonnet 5.5 (claude-sonnet-5-5) for Andr\'as Kornai}
\date{}
\begin{document}
\maketitle

\section{Introduction}
Seychellois Creole, Kreol Seselwa (ISO 639-3 crs; Glottocode sese1246), is a French-based creole and an official national language of Seychelles. According to the search results, Lenstiti Kreol was set up after independence to develop it, and Kreol is the medium of instruction from pre-school to primary 2. This report is short: only a handful of sources were found, and none gives textbooks directly.

\section{Primary Educational Resources}
No textbooks or primers were found with a URL. The academic work below discusses Kreol in education.
\begin{itemize}
\item Zelime and colleagues, work on Kreol Seselwa and its role in education (Academia.edu author page): \url{https://independent.academia.edu/jzelime}
\end{itemize}

\section{Linguistic and Reference Materials}
\begin{itemize}
\item Encyclopedia overview of Seychellois Creole: \url{https://en.wikipedia.org/wiki/Seychellois_Creole}
\item Language-policy overview (in French): \url{https://axl.cefan.ulaval.ca/afrique/seychel.htm}
\item Atlas of Pidgin and Creole Language Structures survey entry: \url{https://apics-online.info/surveys/56}
\item News report on the first monolingual Creole dictionary (about 22,000 words) published by the Seychelles Creole Academy: \url{https://www.lemauricien.com/?p=584344}
\end{itemize}
The search summary also mentioned a Creole--French--English dictionary by Collette Gillieaux, but no URL for it was seen.

\section{Religious and Community Resources}
None found beyond the institutions named above (Lenstiti Kreol, Seychelles Creole Academy), for which no official site was seen.

\section{Digital and Audio Resources}
No audio, software or corpus resources were found.

\section{Curriculum Development Context}
Kreol is already a language of instruction in the early years, so a curriculum can build on existing practice. Extension to upper grades would need the dictionary and orthography resources, whose availability should be verified with the institutions.

\end{document}